\section{Related Work}\label{sec:related}

\paragraph*{Program slicing}
Weiser introduced slicing~\cite{weiser1984program} and showed that programmers use slices when they
debug~\cite{weiser1982programmers}. Ottenstein and Ottenstein~\cite{ottenstein1984pdg}, Ferrante et
al.~\cite{ferrante1987pdg} and, across procedures, Horwitz et
al.~\cite{horwitz1990interprocedural} reduced it to reachability in a dependence graph, and dynamic
slicing restricts a slice to one execution~\cite{korel1988dynamic, agrawal1990dynamic}. Surveys cover the
many variants~\cite{tip1994a, harman2008program}. Frameworks such as WALA slice Python through a separate
front end~\cite{wala}, but they see each gate as a call on one circuit object and cannot tell which qubits
an operation depends on. srcSlice computes static slices from the srcML markup of classical
programs~\cite{alomari2014srcslice}, as our tool does for Qiskit programs, but it has no notion of qubits.
A recent retrospective notes that slices can be defined by almost any semantic method and names quantum
programs as a future direction~\cite{gallagher2025program}. Slicing of probabilistic programs preserves
the distribution of the criterion variables~\cite{hur2014slicing, amtoft2020theory, olmedo2025slicing}.
Preserving $\rho_W$ is stronger, since it preserves the outcome distribution of every measurement of $W$,
and a distribution in one basis does not see phase.

\paragraph*{Dependence in quantum programs}
To our knowledge, Yu and Zhao's QPDG~\cite{yu2025qpdg} is the only prior proposal of dependence-based
slicing for quantum programs. Its dependences are defined in prose, and it leaves an implementation and its
evaluation to future work. Its no-cloning example, \texttt{x(0)} followed by \texttt{cx(0,1)}, is described as
entangling the two qubits, although it produces the product state $\ket{11}$, which is case~(i) of
Proposition~\ref{prop:ent}. The quantum code property graph of Kaul et al.~\cite{kaul2023uniform} adds a
data-flow edge from the control of a CNOT to its target and treats the control as unchanged, the classical
reading that Proposition~\ref{prop:du} refutes. LintQ~\cite{paltenghi2024lintq} orders two operations
exactly when they act on a common qubit, the relation of the QDG, to find bug patterns, aiming by design
at neither full soundness nor full precision. Behler et al.\ extend srcML to OpenQASM, analyze the order of
the operations on each qubit~\cite{behler2024supporting}, and name the slicing of individual qubits as
future work~\cite{behler2025static}.

\paragraph*{Entanglement analysis and light cones}
Static entanglement analyses use type systems~\cite{perdrix2007types, prost2009reasoning, yuan2022twist},
abstract interpretation~\cite{perdrix2008entanglement, yu2021abstract}, the stabilizer
formalism~\cite{honda2015stabilizer} or dataflow rules~\cite{xia2025entanglement}, and the ScaffCC
compiler assumes that interacting qubits may be entangled~\cite{javadiabhari2015scaffcc}, the rule of our
QPDG baseline. They answer whether qubits are correlated, and they are not slicers. Heim et al.\ observe
that dependences between quantum instructions should account for commuting
instructions~\cite{heim2020quantum}, and Qiskit's \texttt{LightCone} pass restricts a circuit to the gates
that can reach an observable~\cite{qiskit2025}. Light-cone arguments bound which gates affect an observable
in the analysis of QAOA~\cite{farhi2014qaoa}, and quantum separation logics reason locally about the qubits
a program touches~\cite{zhou2021qsl, le2022qsl}. Theorem~\ref{thm:sound} states the light-cone fact as a
property of slices of a program.

\paragraph*{Debugging quantum programs}
A recent roadmap finds debugging research for quantum programs in an initial phase, with almost no use of
dependence graphs or data-flow analysis~\cite{ramalho2025road}. The slicing it reviews splits a circuit at
breakpoints the developer places~\cite{metwalli2022debugging, metwalli2024testing}, and Qiskit's slicing
utilities partition a circuit by barriers, depth or gate type~\cite{qiskitaddonutils}. Neither has a
criterion. Bug studies~\cite{zhao2021bugpatterns} and Bugs4Q~\cite{zhao2023bugs4q} catalog real faults,
static checkers such as QChecker~\cite{zhao2023qchecker} detect known bug patterns, mutation
tools~\cite{mendiluze2021muskit, fortunato2022qmutpy} and fault models~\cite{garcia2022testing} supply the
faults of our study, and Ishimoto et al.\ adapt mutation-based fault localization to quantum
programs~\cite{ishimoto2025mbfl}. Slicing restricts where these tools look, although RQ3 shows that the
restriction alone does not improve their ranking.
